\subsection*{Postulates}

\begin{remark*}[Source note]
The following are Euclid's five postulates from Book I of the
\textit{Elements} in the Heath translation
\citep[Book I, Postulates 1--5]{euclid-elements-heath}.
\end{remark*}

\begin{remark*}[Exposition]
Postulates are geometry-specific starting assumptions. They say which
constructions are permitted and which basic geometric facts may be used without
proof. In Euclid's system, these are not conclusions reached from earlier
results; they are part of the foundation from which the theory begins.

Later propositions are proved from the definitions, the postulates, the common
notions, and propositions already established. Thus the postulates are the
rules of play for Euclidean geometry: they determine what can be constructed,
what may be assumed about straight lines and circles, and where the theory's
geometric content enters before proof begins.
\end{remark*}

\begin{axiom}[Euclid Postulate I: Drawing a straight line]
\label{ax:euclid-postulate-draw-straight-line}
Let it have been postulated to draw a straight line from any point to any point.
\end{axiom}
\begin{remark*}[Standard quantified statement]
This is a classical Euclidean source-text clause. In this chapter it is treated
as a named primitive statement rather than expanded into a modern first-order
formalization.
\end{remark*}

\begin{lra-not-visible}
\begin{dependencies}
\begin{itemize}
  \item \hyperref[def:euclid-point]{Point}
  \item \hyperref[def:euclid-straight-line]{Straight line}
\end{itemize}
\end{dependencies}
\end{lra-not-visible}

\begin{axiom}[Euclid Postulate II: Producing a finite straight line]
\label{ax:euclid-postulate-produce-straight-line}
Let it have been postulated to produce a finite straight line continuously in a
straight line.
\end{axiom}
\begin{remark*}[Standard quantified statement]
This is a classical Euclidean source-text clause. In this chapter it is treated
as a named primitive statement rather than expanded into a modern first-order
formalization.
\end{remark*}

\begin{lra-not-visible}
\begin{dependencies}
\begin{itemize}
  \item \hyperref[def:euclid-straight-line]{Straight line}
\end{itemize}
\end{dependencies}
\end{lra-not-visible}

\begin{axiom}[Euclid Postulate III: Describing a circle]
\label{ax:euclid-postulate-describe-circle}
Let it have been postulated to describe a circle with any center and distance.
\end{axiom}
\begin{remark*}[Standard quantified statement]
This is a classical Euclidean source-text clause. In this chapter it is treated
as a named primitive statement rather than expanded into a modern first-order
formalization.
\end{remark*}

\begin{lra-not-visible}
\begin{dependencies}
\begin{itemize}
  \item \hyperref[def:euclid-center-of-a-circle]{Center of a circle}
  \item \hyperref[def:euclid-circle]{Circle}
  \item \hyperref[def:euclid-point]{Point}
\end{itemize}
\end{dependencies}
\end{lra-not-visible}

\begin{axiom}[Euclid Postulate IV: Equality of right angles]
\label{ax:euclid-postulate-right-angles-equal}
Let it have been postulated that all right angles are equal to one another.
\end{axiom}
\begin{remark*}[Standard quantified statement]
This is a classical Euclidean source-text clause. In this chapter it is treated
as a named primitive statement rather than expanded into a modern first-order
formalization.
\end{remark*}

\begin{lra-not-visible}
\begin{dependencies}
\begin{itemize}
  \item \hyperref[def:euclid-right-angle-and-perpendicular]{Right angle and perpendicular}
\end{itemize}
\end{dependencies}
\end{lra-not-visible}

\begin{axiom}[Euclid Postulate V: Parallel postulate]
\label{ax:euclid-postulate-parallel}
Let it have been postulated that, if a straight line falling on two straight
lines makes the interior angles on the same side less than two right angles,
the two straight lines, if produced indefinitely, meet on that side on which are
the angles less than the two right angles.
\end{axiom}
\begin{remark*}[Standard quantified statement]
This is a classical Euclidean source-text clause. In this chapter it is treated
as a named primitive statement rather than expanded into a modern first-order
formalization.
\end{remark*}

\begin{lra-not-visible}
\begin{dependencies}
\begin{itemize}
  \item \hyperref[def:euclid-plane-angle]{Plane angle}
  \item \hyperref[def:euclid-right-angle-and-perpendicular]{Right angle and perpendicular}
  \item \hyperref[def:euclid-straight-line]{Straight line}
\end{itemize}
\end{dependencies}
\end{lra-not-visible}
